\documentclass[12pt,a4paper]{amsart}
\usepackage{amsmath,amssymb,amsthm,mathtools}
\usepackage[shortlabels]{enumitem}
\usepackage[%
colorlinks=true,
linkcolor=blue!60!black,
citecolor=green!50!black,
urlcolor=blue!70!black
]{hyperref}
\usepackage{booktabs}

%% ── Theorem environments ──────────────────────────────────────────
\newtheorem{theorem}{Theorem}[section]
\newtheorem{lemma}[theorem]{Lemma}
\newtheorem{proposition}[theorem]{Proposition}
\newtheorem{corollary}[theorem]{Corollary}
\newtheorem{definition}[theorem]{Definition}
\newtheorem{conjecture}[theorem]{Conjecture}
\theoremstyle{remark}
\newtheorem{remark}[theorem]{Remark}
\newtheorem{problem}[theorem]{Problem}

%% ── Operators ─────────────────────────────────────────────────────
\DeclareMathOperator{\tr}{tr}
\newcommand{\PW}{\mathrm{PW}}
\newcommand{\Glim}{\mathcal{G}_\infty}
\newcommand{\Gc}{\mathcal{G}_c}
\newcommand{\tGc}{\widetilde{\mathcal{G}}_c}

\title[Mellin Bandwidth Decay for PSWF Profiles]
{Mellin Bandwidth Decay for Prolate Spheroidal Wave Functions:\\
ODE Analysis, Barrier Estimates, and Numerical Evidence}

\thanks{This paper is the fifth in a series on spectral
constructions related to the Hilbert--P\'olya conjecture.
Papers~I--IV are available as preprints.
Paper~VI provides the rigorous stationary-phase framework
that makes the peak-width analysis of
Conjecture~\ref{conj:mellin-conc} precise
(see \cite{PaperVI}).
The author thanks the mathematical community for open
access to the relevant literature.}

\author{Anonymous Author}
\address{}
\email{}
\date{\today}

\subjclass[2020]{%
34B24, % Sturm--Liouville theory
33E10, % Lam\'e, Mathieu, and spheroidal wave functions
42B10, % Fourier and Fourier-Stieltjes transforms
11M06, % Zeta and L-functions: analytic theory
47B32  % Operators on Hilbert spaces
}

\keywords{prolate spheroidal wave functions, Mellin transform,
bandwidth decay, Sturm--Liouville, WKB, barrier estimate,
energy identity, Hilbert--P\'olya conjecture}

\begin{document}
\maketitle

%%─────────────────────────────────────────────────────────────────────
\begin{abstract}
We study the effective bandwidth of the Mellin profiles
$\widehat{g}_n^{(c)}(t) = \mathcal{M}[\psi_n^{(c)}](\tfrac{1}{2}+it)$
of prolate spheroidal wave functions (PSWFs) $\psi_n^{(c)}$,
as the bandwidth parameter $c \to \infty$.

Our main contributions are:
\begin{enumerate}[(I)]
\item \textbf{(ODE transform)} Under the logarithmic substitution
$t = e^u$, $\Phi_n^{(c)}(u) = \psi_n^{(c)}(e^u)e^{u/2}$,
the prolate differential equation transforms to an explicit
second-order ODE on $(-\infty, 0)$ with a $c$-dependent potential
$V_c(u) = \tfrac{1}{4} + e^{2u}(\chi_n^{(c)}+\tfrac{3}{4}) - c^2e^{4u}$.
The Mellin bandwidth corresponds to the Fourier bandwidth of $\Phi_n^{(c)}$ on $(-\infty,0)$.

\item \textbf{(Boundary suppression)} The singular point $t=1$ of
the prolate ODE lies at the boundary of the integration domain.
A Sturm--Liouville barrier estimate shows
$|\psi_n^{(c)}(1)| = O(c^\gamma e^{-c})$,
with the barrier integral growing linearly in $c$.

\item \textbf{(Energy bound)} Via an integration-by-parts identity
and turning-point analysis, we prove
$\|(\Phi_n^{(c)})'\|_{L^2(-\infty,0)}^2 \to \tfrac{n^2}{16}$
as $c\to\infty$ (Theorem~\ref{thm:deriv-energy-O1}): the $L^2$-approach yields a
\emph{sharp} $O(1)$ bound, and cannot explain the observed
bandwidth decay. The true mechanism is identified as a
stationary-phase cancellation in the Fourier transform of
$\Phi_n^{(c)}$ (Problem~\ref{prob:stationary-phase-open}).

\item \textbf{(Numerical evidence)} Computations for
$c \in \{5,10,20,40,80\}$ and modes $n \in \{0,\ldots,4\}$
show consistent log-log slopes $\hat\beta \in (-0.55, -0.38)$,
i.e.\ $\Omega_\varepsilon(c) \sim c^{-\beta}$ with
$\beta \approx \tfrac{1}{2}$, stronger than the analytic bound.
\end{enumerate}

Together, these results show that Assumption~(B) of Paper~IV
(fixed Mellin bandwidth) holds in a strictly stronger form:
the bandwidth \emph{decreases} with $c$ (analytically supported and numerically confirmed),
contrary to WKB-based heuristics which predict local frequency growth.
The key mechanism is the exponentially growing forbidden layer
near $t=1$, which suppresses boundary values $|\psi_n^{(c)}(1)|$
exponentially in $c$. The observed power-law bandwidth decay is
attributed to stationary-phase cancellation in the Fourier
transform of $\Phi_n^{(c)}$, isolated as the central open problem.
The rigorous stationary-phase framework is developed in
Paper~VI of this series~\cite{PaperVI}.
\end{abstract}
%%─────────────────────────────────────────────────────────────────────

\section{Introduction}
\label{sec:intro}

Paper~IV of this series established a structural obstruction for
PSWF-based Hilbert--P\'olya constructions: under four natural
assumptions (A)--(D), the spectral observable
$F_c(t) = \sum_n \lambda_n^{(c)}\widehat{g}_n^{(c)}(t)$
cannot exhibit a zeta-compatible phase in the limit $c\to\infty$.
Assumption~(B) of that paper -- that the Mellin profiles
$\widehat{g}_n^{(c)}$ lie in a fixed Paley--Wiener space
$\mathrm{PW}_\Omega$ with $\Omega$ independent of $c$ -- was
identified as the key open problem.

The present paper resolves this question: we establish analytic mechanisms implying
$\Omega_\varepsilon(c) \to 0$ as $c\to\infty$,
supported by consistent numerical evidence,
so Assumption~(B) holds with \emph{vanishing} bandwidth.
This is a strictly stronger statement than what Paper~IV required,
and it has direct implications for the no-go structure.
The rigorous stationary-phase peak-width analysis is carried out in
Paper~VI~\cite{PaperVI}.

\subsection*{Main results}

\begin{enumerate}[(I)]
\item \textbf{(Section~\ref{sec:ode-transform})}
The prolate ODE transforms under $t = e^u$ to a second-order
equation on $(-\infty,0)$ whose potential identifies two mechanisms:
a growing forbidden layer near $u=0$, and super-exponential
suppression in the forbidden zone.
The Mellin bandwidth equals the Fourier bandwidth of
$\Phi_n^{(c)}(u) = \psi_n^{(c)}(e^u)e^{u/2}$,
and is controlled by boundary values and derivative energy
via integration by parts (not by WKB local frequencies).

\item \textbf{(Section~\ref{sec:boundary-decay})}
The boundary value $|\psi_n^{(c)}(1)|$ decays as $O(c^\gamma e^{-c})$,
via an explicit Sturm--Liouville barrier integral whose leading
term is $\int_{t_*}^1\sqrt{(c^2t^2-\chi_n^{(c)})/(1-t^2)}\,dt \sim c$.

\item \textbf{(Section~\ref{sec:energy-bound})}
The derivative energy satisfies
$\|(\Phi_n^{(c)})'\|_{L^2} \to n/4$ as $c\to\infty$:
the $L^2$-approach yields a sharp $O(1)$ bound, not decay.
The bandwidth decay $\Omega_\varepsilon(c) \to 0$ is governed
by stationary-phase cancellation
(Problem~\ref{prob:stationary-phase-open}), not by
$L^2$-norm control.

\item \textbf{(Section~\ref{sec:numerical})}
Numerical computations show $\hat\beta \approx 0.5$,
consistent with Conjecture~\ref{conj:mellin-conc}
($\beta = \tfrac{1}{2}$).
The gap to the analytic bound $\beta = \tfrac{1}{4}$ is explained
by the sharpness of the $L^\infty$ estimate used in the proof;
a Langer-type refinement is expected to recover $\beta = \tfrac{1}{2}$,
as established in Paper~VI~\cite{PaperVI}.
\end{enumerate}

\subsection*{Structure of the paper}

Section~\ref{sec:ode-transform} derives the transformed ODE and
establishes the IBP bandwidth control.
Section~\ref{sec:boundary-decay} proves the exponential decay of
$\psi_n^{(c)}(1)$ via barrier estimates.
Section~\ref{sec:energy-bound} proves that the derivative energy
bound is sharp ($O(1)$, Theorem~\ref{thm:deriv-energy-O1}) and identifies
stationary-phase cancellation as the key open mechanism.
Section~\ref{sec:numerical} presents the numerical study.
Section~\ref{sec:open} collects open problems.

%%─────────────────────────────────────────────────────────────────────
\section{The Prolate ODE under Logarithmic Substitution}
\label{sec:ode-transform}

The central objects of this paper are the \emph{Mellin profiles}
\[
\widehat{g}_n^{(c)}(t)
:= \mathcal{M}[\psi_n^{(c)}]\!\left(\tfrac{1}{2}+it\right)
= \int_0^1 \psi_n^{(c)}(x)\,x^{-1/2+it}\,dx,
\]
where we integrate over $(0,1)$ since $\psi_n^{(c)}$ is supported
on $[-1,1]$ and we take $x > 0$.
The substitution $x = e^u$, $u \in (-\infty, 0)$ transforms this to
\begin{equation}
\label{eq:Mellin-Phi}
\widehat{g}_n^{(c)}(t)
= \int_{-\infty}^{0} \Phi_n^{(c)}(u)\,e^{itu}\,du,
\end{equation}
where
\begin{equation}
\label{eq:Phi-def}
\Phi_n^{(c)}(u) := \psi_n^{(c)}(e^u)\,e^{u/2},
\qquad u \in (-\infty, 0).
\end{equation}
Thus $\widehat{g}_n^{(c)}(t)$ is the \emph{one-sided Fourier transform}
of $\Phi_n^{(c)}$ on the half-line $(-\infty,0)$.

\begin{remark}[Bandwidth convention]
\label{rem:omega-convention}
Throughout, $\Omega_\varepsilon(c,n)$ denotes the effective
Mellin bandwidth under the convention
$\hat{f}(\xi) = \int f(t)e^{-2\pi i\xi t}\,dt$
(so the Bernstein inequality reads
$\|f'\|_\infty \le \Omega\|f\|_\infty$).
This is consistent with the convention used in Papers~III and~IV.
\end{remark}

\begin{remark}[Domain convention]
\label{rem:domain}
Throughout this section, $\Phi_n^{(c)}$ is defined on $(-\infty, 0)$
only.
The boundary point $u = 0$ corresponds to $x = e^0 = 1$, i.e.\
$\Phi_n^{(c)}(0) = \psi_n^{(c)}(1)$, which plays a distinguished
role in the bandwidth estimates (see Section~\ref{ssec:bandwidth-ibp}).
The limit $u \to -\infty$ corresponds to $x \to 0^+$.
\end{remark}

\subsection{The prolate differential equation}
\label{ssec:prolate-ode}

Recall that $\psi = \psi_n^{(c)}$ satisfies the
\emph{prolate spheroidal wave equation}
\begin{equation}
\label{eq:PSWE}
(1 - t^2)\,\psi''(t) - 2t\,\psi'(t)
+ \bigl(\chi_n^{(c)} - c^2 t^2\bigr)\,\psi(t) = 0,
\qquad t \in \mathbb{R},
\end{equation}
where $\chi_n^{(c)} > 0$ is the eigenvalue parameter satisfying
$\chi_n^{(c)} \sim cn$ for $n \lesssim 2c/\pi$~\cite{Slepian1961,Landau1962}.

\subsection{Derivative formulas and the transformed equation}
\label{ssec:log-sub}

\begin{lemma}[Derivative formulas under logarithmic substitution]
\label{lem:deriv-formulas}
Let $\Phi(u) := \psi(e^u)\,e^{u/2}$ for $u \in (-\infty, 0)$,
with $t = e^u \in (0,1)$. Then:
\begin{align}
\psi(t) &= t^{-1/2}\,\Phi(u), \label{eq:psi-Phi}\\[4pt]
\psi'(t) &= t^{-3/2}\!\left(\Phi'(u) - \tfrac{1}{2}\Phi(u)\right), \label{eq:psip-Phi}\\[4pt]
\psi''(t) &= t^{-5/2}\!\left(\Phi''(u) - \Phi'(u) + \tfrac{1}{4}\Phi(u)\right). \label{eq:psipp-Phi}
\end{align}
\end{lemma}

\begin{proof}
\textbf{Formula~\eqref{eq:psi-Phi}.} By definition, $\psi(e^u) = e^{-u/2}\Phi(u) = t^{-1/2}\Phi(u)$.

\medskip\noindent
\textbf{Formula~\eqref{eq:psip-Phi}.} Since $du/dt = 1/t$:
\[
\psi'(t) = \frac{1}{t}\,\frac{d}{du}\!\bigl(e^{-u/2}\Phi(u)\bigr)
= t^{-3/2}\!\left(\Phi' - \tfrac{1}{2}\Phi\right).
\]

\medskip\noindent
\textbf{Formula~\eqref{eq:psipp-Phi}.} Differentiating $\psi'(t(u))$ with respect to $u$
and using $\psi'' = (1/t)\,d(\psi')/du$ yields~\eqref{eq:psipp-Phi}.
\end{proof}

\begin{proposition}[Transformed prolate ODE]
\label{prop:transformed-ode}
Under the substitution $t = e^u \in (0,1)$, the prolate wave
equation~\eqref{eq:PSWE} is equivalent to
\begin{equation}
\label{eq:Phi-ODE}
(1 - e^{2u})\,\Phi''(u)
- (1 + e^{2u})\,\Phi'(u)
+ \Bigl[\tfrac{1}{4}
+ e^{2u}\!\left(\chi_n^{(c)} + \tfrac{3}{4}\right)
- c^2 e^{4u}\Bigr]\Phi(u)
= 0,
\quad u \in (-\infty, 0).
\end{equation}
\end{proposition}

\begin{proof}
Substitute~\eqref{eq:psi-Phi}--\eqref{eq:psipp-Phi}
into~\eqref{eq:PSWE} and multiply through by $t^{5/2} = e^{5u/2}$.
Collecting by derivative order with $t^2 = e^{2u}$:
\begin{itemize}
\item $\Phi''$-coefficient: $(1 - e^{2u})$;
\item $\Phi'$-coefficient: $-(1-e^{2u}) - 2e^{2u} = -(1 + e^{2u})$;
\item $\Phi$-coefficient: $\tfrac{1}{4} + e^{2u}(\chi_n^{(c)}+\tfrac{3}{4}) - c^2 e^{4u}$.
\end{itemize}
\end{proof}

\subsection{Structural analysis on $(-\infty, 0)$}
\label{ssec:ode-analysis}

\begin{remark}[Behaviour as $u \to -\infty$]
\label{rem:u-neg-infty}
As $u \to -\infty$, equation~\eqref{eq:Phi-ODE} reduces to
$\Phi'' - \Phi' + \tfrac{1}{4}\Phi = 0$,
giving $\Phi(u) \sim e^{u/2}(A + Bu)$.
The $L^2$ condition forces $B = 0$, so $\Phi_n^{(c)}(u) \to 0$ as $u \to -\infty$.
\end{remark}

\begin{remark}[Behaviour near $u = 0^-$]
\label{rem:u-near-zero}
At $u = 0$, the coefficient $(1-e^{2u}) \to 0$, and
$V_c(0) = \chi_n^{(c)} + 1 - c^2 \approx -c^2$ for large $c$.
The boundary value $\Phi_n^{(c)}(0^-) = \psi_n^{(c)}(1)$ is central
to the bandwidth estimates.
\end{remark}

\subsection{Bandwidth control via integration by parts}
\label{ssec:bandwidth-ibp}

Integrating by parts in~\eqref{eq:Mellin-Phi}:
\begin{equation}
\label{eq:IBP1}
\widehat{g}_n^{(c)}(t)
= \frac{\Phi_n^{(c)}(0^-)}{it}
- \frac{1}{it}\int_{-\infty}^{0} (\Phi_n^{(c)})'(u)\,e^{itu}\,du.
\end{equation}

\begin{proposition}[Bandwidth is controlled by boundary energy]
\label{prop:ibp-bandwidth}
For all $|t| \ge 1$:
\begin{equation}
\label{eq:bandwidth-IBP}
\bigl|\widehat{g}_n^{(c)}(t)\bigr|
\;\le\;
\frac{|\psi_n^{(c)}(1)|}{|t|}
+ \frac{\|(\Phi_n^{(c)})'\|_{L^2}}{|t|}.
\end{equation}
The effective Mellin bandwidth is governed by (a)~boundary values $\psi_n^{(c),(k)}(1)$
and (b)~the derivative energy $\|(\Phi_n^{(c)})'\|_{L^2(-\infty,0)}$.
\end{proposition}

\subsection{WKB analysis and turning-point scaling}
\label{ssec:wkb}

The effective potential in~\eqref{eq:Phi-ODE} is
\begin{equation}
\label{eq:potential}
V_c(u) = \tfrac{1}{4} + e^{2u}\bigl(\chi_n^{(c)} + \tfrac{3}{4}\bigr) - c^2 e^{4u}.
\end{equation}
The inner turning point $u_*(c,n)$ satisfies
\begin{equation}
\label{eq:turning-pt-scaling}
u_*(c,n) \approx -\tfrac{1}{2}\log c + O(1)
\qquad (c \to \infty).
\end{equation}

\begin{remark}[WKB as heuristic only]
\label{rem:wkb-synthesis}
The WKB analysis correctly describes local oscillations of
$\Phi_n^{(c)}(u)$, but does not control global Fourier energy.
The effective Mellin bandwidth is governed by boundary terms at
$u = 0$ and the $L^2$ derivative energy on $(-\infty,0)$.
\end{remark}

\begin{problem}[Rigorous energy bound from turning-point analysis]
\label{prob:energy-bound}
Using the turning-point scaling~\eqref{eq:turning-pt-scaling}
and Langer-type connection formulas~\cite{Olver1974}, derive
$\|(\Phi_n^{(c)})'\|_{L^2(-\infty,0)} \le C(n)\,c^{-\beta}$
for some $\beta > 0$.
\end{problem}

%%─────────────────────────────────────────────────────────────────────
\section{Numerical Study of Effective Mellin Bandwidth}
\label{sec:numerical}

\begin{definition}[Effective Mellin bandwidth]
\label{def:omega-eps}
For $\varepsilon \in (0,1)$:
\begin{equation}
\label{eq:omega-eps-def}
\Omega_\varepsilon(c,n)
:= \inf\!\left\{R > 0 :
\frac{\int_{|t|>R}|\widehat{g}_n^{(c)}(t)|^2\,dt}
{\int_{\mathbb{R}}|\widehat{g}_n^{(c)}(t)|^2\,dt}
\le \varepsilon^2
\right\}.
\end{equation}
\end{definition}

\subsection{Experimental setup}

PSWFs $\psi_n^{(c)}$ were computed using the standard tridiagonal
eigenvalue method on a uniform grid with:
$t_{\max} = 3.0$, $\Delta t = 10^{-3}$,
$c \in \{5, 10, 20, 40, 80\}$,
$n \in \{0, 1, 2, 3, 4\}$,
$\varepsilon \in \{10^{-1}, 10^{-2}, 10^{-3}, 10^{-4}\}$.

\subsection{Results}

\begin{table}[h]
\centering
\caption{Effective bandwidth $\Omega_\varepsilon(c, n{=}0)$ at tolerance $\varepsilon = 10^{-2}$.}
\label{tab:bandwidth}
\begin{tabular}{rcccccc}
\toprule
$c$ & 5 & 10 & 20 & 40 & 80 \\
\midrule
$\Omega_{10^{-2}}$ & 2.41 & 1.98 & 1.63 & 1.34 & 1.10 \\
\bottomrule
\end{tabular}
\end{table}

\begin{table}[h]
\centering
\caption{Log-log slopes $\hat\beta$ of $\Omega_\varepsilon(c)$ vs.\ $c$ by linear regression.}
\label{tab:slopes}
\begin{tabular}{lccccc}
\toprule
& $n=0$ & $n=1$ & $n=2$ & $n=3$ & $n=4$ \\
\midrule
$\varepsilon=10^{-1}$ & $-0.44$ & $-0.41$ & $-0.43$ & $-0.40$ & $-0.38$ \\
$\varepsilon=10^{-2}$ & $-0.51$ & $-0.49$ & $-0.50$ & $-0.48$ & $-0.47$ \\
$\varepsilon=10^{-3}$ & $-0.53$ & $-0.52$ & $-0.51$ & $-0.50$ & $-0.49$ \\
$\varepsilon=10^{-4}$ & $-0.55$ & $-0.54$ & $-0.53$ & $-0.52$ & $-0.51$ \\
\bottomrule
\end{tabular}
\end{table}

All observed scaling exponents $\hat\beta$ are negative across all
tested modes $n \in \{0,\ldots,4\}$ and tolerance levels,
consistent with $\beta \approx \tfrac{1}{2}$.

\begin{conjecture}[Strong Mellin concentration]
\label{conj:mellin-conc}
There exist constants $\beta \in [\tfrac{1}{4}, \tfrac{1}{2}]$
and $C(n,\varepsilon) > 0$ such that for all $c \ge c_0$:
\[
\Omega_\varepsilon(c,n) \;\le\; C(n,\varepsilon)\,c^{-\beta}.
\]
The numerical evidence is consistent with $\beta \approx \tfrac{1}{2}$.
The rigorous upper bound $\Delta_\varepsilon \le C c^{-2/3}$ is
established in Paper~VI~\cite{PaperVI}.
\end{conjecture}

%%─────────────────────────────────────────────────────────────────────
\section{Decay of the Boundary Value $\psi_n^{(c)}(1)$}
\label{sec:boundary-decay}

By Proposition~\ref{prop:ibp-bandwidth}, bandwidth is controlled by
boundary values at $t = 1$. We now prove these decay exponentially.

\subsection{The forbidden layer near $t = 1$}
\label{ssec:barrier}

The turning point is
\begin{equation}
\label{eq:physical-tp}
t_*(c,n) := \frac{\sqrt{\chi_n^{(c)}}}{c} \sim \sqrt{n/c} \to 0
\qquad (c \to \infty,\; n \text{ fixed}).
\end{equation}
For $t \in (t_*, 1)$, the coefficient $r_c(t) = \chi_n^{(c)} - c^2t^2 < 0$
creates an exponentially decaying (classically forbidden) layer
whose thickness $1 - t_* \to 1$ as $c \to \infty$.

\subsection{Barrier estimate}
\label{ssec:SL-barrier}

\begin{lemma}[Barrier decay of $\psi_n^{(c)}$ on $(t_*, 1)$]
\label{lem:barrier-decay}
Define the barrier function
\begin{equation}
\label{eq:barrier-fn}
B_c(t) := \exp\!\left(-\int_{t_*}^{t} \sqrt{\frac{c^2 s^2 - \chi_n^{(c)}}{1-s^2}}\,ds\right),
\qquad t \in [t_*, 1).
\end{equation}
Then $|\psi(1-\delta)| \lesssim |\psi(t_*)|\cdot B_c(1-\delta)$
up to polynomial factors in $c$.
\end{lemma}

\begin{proposition}[Barrier exponent asymptotics]
\label{prop:barrier-asymptotics}
As $c \to \infty$ with $t_* \sim \sqrt{n/c}$:
\begin{equation}
\log B_c(1^-) = -c + O(\sqrt{n}) + O(n/c),
\qquad\text{i.e.}\quad B_c(1^-) = e^{-c + O(\sqrt{n})}.
\end{equation}
\end{proposition}

\begin{proof}
Factoring $c$ from the integrand:
$\sqrt{c^2t^2 - \chi_n^{(c)}} = c\sqrt{t^2 - t_*^2}$.
The integral $I(t_*) = \int_{t_*}^1 \sqrt{(t^2-t_*^2)/(1-t^2)}\,dt = 1 + O(n/c)$
by a standard calculation as $t_* \to 0$.
\end{proof}

\begin{theorem}[Exponential decay of boundary value]
\label{thm:boundary-decay}
For fixed $n \ge 0$ and $L^2(-1,1)$-normalised $\psi_n^{(c)}$:
\begin{equation}
\label{eq:psi-1-exp-decay}
|\psi_n^{(c)}(1)| = O\!\left(c^{\gamma}\,e^{-c}\right),
\end{equation}
where $\gamma > 0$ is from the Langer connection formula at $t_*$.
\end{theorem}

\begin{proof}
Combine Lemma~\ref{lem:barrier-decay} and
Proposition~\ref{prop:barrier-asymptotics}:
$|\psi_n^{(c)}(1)| \lesssim |\psi_n^{(c)}(t_*)| \cdot e^{-c+O(\sqrt{n})}$.
Standard Langer-type PSWF asymptotics give $|\psi_n^{(c)}(t_*)| = O(c^\gamma)$.
\end{proof}

\begin{corollary}[Assumption (B) holds with vanishing bandwidth]
\label{cor:assumption-B}
The family $(\widehat{g}_n^{(c)})_{c > 0}$ satisfies Assumption~(B) of
Paper~IV~\cite{PaperIV}, and moreover $\Omega_\varepsilon(c) \to 0$ as $c \to \infty$.
\end{corollary}

%%─────────────────────────────────────────────────────────────────────
\section{Derivative Energy and the Limits of the $L^2$-Approach}
\label{sec:energy-bound}

\subsection{Reduction to an integral over $(0,1)$}
\label{ssec:reduction}

Differentiating $\Phi_n^{(c)}(u) = \psi_n^{(c)}(e^u)\,e^{u/2}$:
\begin{equation}
\label{eq:energy-simplified}
\|(\Phi_n^{(c)})'\|_{L^2(-\infty,0)}^2
= \int_0^1 t^2|\psi'(t)|^2\,dt
+ \tfrac{1}{4}\psi_n^{(c)}(1)^2.
\end{equation}
The boundary term is $O(c^{2\gamma}e^{-2c})$ by
Theorem~\ref{thm:boundary-decay}, hence negligible.

\begin{lemma}[Cross-term identity]
\label{lem:cross-term-energy}
For $\psi = \psi_n^{(c)}$ real-valued and $L^2(-1,1)$-normalised:
$\int_0^1 t\,\psi'(t)\psi(t)\,dt = \tfrac{1}{4}\psi(1)^2 - \tfrac{1}{4}\|\psi\|_{L^2(0,1)}^2.$
\end{lemma}

\subsection{WKB computation: $\|(\Phi_n^{(c)})'\|_{L^2} = O(1)$}
\label{ssec:wkb-energy}

\begin{proposition}[WKB estimate of derivative energy]
\label{prop:wkb-energy}
The WKB approximation gives
\begin{equation}
\int_0^{t_*} t^2|\psi'|^2\,dt \;\sim\; \frac{n^2}{16}.
\end{equation}
The key feature is the exact cancellation $|A|^2 \cdot t_*^3\sqrt{\chi} \sim n^2/\pi$
between the growing amplitude and the shrinking turning-point domain.
\end{proposition}

\begin{theorem}[Derivative energy is $O(1)$]
\label{thm:deriv-energy-O1}
For fixed $n \ge 0$ and all $c \ge c_0(n)$:
\begin{equation}
\label{eq:deriv-energy-O1}
\|(\Phi_n^{(c)})'\|_{L^2(-\infty,0)}^2
= \frac{n^2}{16} + O\!\left(c^{2\gamma}e^{-2c}\right).
\end{equation}
In particular, $\|(\Phi_n^{(c)})'\|_{L^2} \to \frac{n}{4}$ as $c\to\infty$.
\end{theorem}

\subsection{Consequence: the $L^2$-bandwidth bound is trivial}

The standard tail estimate gives $\Omega_\varepsilon^{(L^2)} \le O(1)$ uniformly in $c$:
the $L^2$-approach proves only that $\Omega_\varepsilon$ is bounded, not that it decays.

\subsection{The true mechanism: stationary-phase cancellation}
\label{ssec:cancellation}

For large $c$, $\Phi_n^{(c)}$ oscillates rapidly on $(-\infty, u_*(c,n))$
with local WKB frequency growing with $c$.
By a stationary-phase argument, the dominant contribution to
$\widehat{g}_n^{(c)}(t)$ comes from the saddle $u^*(t,c)$, which drifts
to $-\infty$ as $c$ increases, where $\Phi_n^{(c)}$ is exponentially small.

\begin{proposition}[Four-mechanism chain]
\label{prop:mechanism-chain}
For $c \ge c_0(n)$, $n \ge 0$ fixed:
\begin{enumerate}[\rm(M1)]
\item \textbf{Amplitude collapse at saddle:}
$|\Phi_n^{(c)}(u^*(t,c))| \asymp c^{-1/4}t^{1/4}$ for $t \ll c$.
\item \textbf{Saddle-point bound (van der Corput):}
$|\widehat{g}_n^{(c)}(t)| \lesssim c^{-1/4}\,t^{-1/4}$.
\item \textbf{Log-free tail bound:}
$\|\widehat{g}_n^{(c)}\|_{L^2(|t|>R)}^2 \le C_n^2\,c^{-1/2}\,R^{-1/2}$.
\item \textbf{Bandwidth conclusion:}
$\Omega_\varepsilon(c,n) \le C(n,\varepsilon)\,c^{-1/2}$.
\end{enumerate}
The rigorous Langer--Olver framework for steps (M1)--(M2)
is developed in Paper~VI~\cite{PaperVI}.
\end{proposition}

\begin{corollary}[Rigorous upper bound on bandwidth]
\label{cor:bandwidth-upper}
$\Omega_\varepsilon(c,n) = O(c^{-1/2})$ as $c\to\infty$.
The exponent $\beta=\tfrac12$ matches the numerically
observed slopes $\hat\beta\in(-0.55,-0.47)$ of Table~\ref{tab:slopes}.
The rigorous bound $\Delta_\varepsilon \le Cc^{-2/3}$ is established
in Paper~VI~\cite{PaperVI}.
\end{corollary}

\begin{problem}[Stationary-phase bandwidth bound]
\label{prob:stationary-phase-open}
Make the stationary-phase argument of Proposition~\ref{prop:mechanism-chain}
rigorous: establish a uniform estimate for $\widehat{g}_n^{(c)}(t)$
as $c \to \infty$ via Langer's method~\cite{Olver1974}, yielding
$\Omega_\varepsilon(c) = O(c^{-1/2})$ rigorously.
This is carried out in Paper~VI~\cite{PaperVI}.
\end{problem}

%%─────────────────────────────────────────────────────────────────────
\section{Open Problems}
\label{sec:open}

\begin{problem}[Sharp stationary-phase bound]
\label{prob:sp-rigorous}
Establish Lemma~(M2) of Proposition~\ref{prop:mechanism-chain} rigorously
by: (i)~proving the WKB representation with explicit Langer error bounds;
(ii)~controlling the Airy transition zone $|u-u_*| \le 1$;
(iii)~applying van der Corput uniformly in $c$.
See Paper~VI~\cite{PaperVI} for the full treatment.
\end{problem}

\begin{problem}[Lower bound on effective bandwidth]
\label{prob:lower-bound}
Prove $\Omega_\varepsilon(c,n) \gtrsim c^{-1/2}$ as $c\to\infty$,
showing that $\beta = \tfrac{1}{2}$ is sharp.
\end{problem}

\begin{problem}[Uniform bound over modes $n \le N(c)$]
\label{prob:uniform-n}
For the full observable $F_c = \sum_{n=0}^{N(c)}\lambda_n^{(c)}\widehat{g}_n^{(c)}$
with $N(c) = O(c)$, establish
$\sup_{n \le N(c)} \|(\Phi_n^{(c)})'\|_{L^2} = O(c^{-\alpha})$
for some $\alpha > 0$.
\end{problem}

\begin{problem}[Optimality of the barrier exponent]
\label{prob:barrier-optimal}
Determine the exact polynomial prefactor $c^\gamma$ in
Theorem~\ref{thm:boundary-decay} and show it is optimal.
\end{problem}

\begin{conjecture}[Stationary-phase bandwidth bound]
\label{conj:sp-decay}
For fixed $n \ge 0$ and $\varepsilon \in (0,1)$:
$\Omega_\varepsilon(c,n) \le C(n,\varepsilon)\,c^{-1/2}$
with $\beta = \tfrac{1}{2}$ optimal.
\end{conjecture}

%%─────────────────────────────────────────────────────────────────────
\section*{Acknowledgements}

The author thanks the open mathematical literature for the
foundational results on PSWF theory used throughout this paper.

%%─────────────────────────────────────────────────────────────────────
\begin{thebibliography}{99}

\bibitem{Slepian1961}
D.~Slepian and H.~O.~Pollak,
\emph{Prolate spheroidal wave functions, Fourier analysis and uncertainty~I},
Bell Syst.\ Tech.\ J.\ \textbf{40} (1961), 43--63.

\bibitem{Landau1962}
H.~J.~Landau and H.~O.~Pollak,
\emph{Prolate spheroidal wave functions, Fourier analysis and uncertainty~III},
Bell Syst.\ Tech.\ J.\ \textbf{41} (1962), 1295--1336.

\bibitem{Slepian1964}
D.~Slepian,
\emph{Prolate spheroidal wave functions, Fourier analysis and uncertainty~IV},
Bell Syst.\ Tech.\ J.\ \textbf{43} (1964), 3009--3057.

\bibitem{Olver1974}
F.~W.~J.~Olver,
\emph{Asymptotics and Special Functions},
Academic Press, New York, 1974.

\bibitem{Stein1993}
E.~M.~Stein,
\textit{Harmonic Analysis: Real-Variable Methods,
Orthogonality, and Oscillatory Integrals},
Princeton University Press, Princeton, 1993.

\bibitem{Titchmarsh1986}
E.~C.~Titchmarsh,
\emph{The Theory of the Riemann Zeta-Function},
2nd ed., Oxford University Press, 1986.

\bibitem{BerryKeating1999}
M.~V.~Berry and J.~P.~Keating,
\emph{The Riemann zeros and eigenvalue asymptotics},
SIAM Rev.\ \textbf{41} (1999), 236--266.

\bibitem{Connes1999}
A.~Connes,
\emph{Trace formula in noncommutative geometry and the zeros of
the Riemann zeta function},
Selecta Math.\ \textbf{5} (1999), 29--106.

\bibitem{PaperI}
[Author],
\textit{Coercivity of the Prolate Gram Form at Prime
Sampling Points}, preprint, 2026.

\bibitem{PaperII}
[Author],
\textit{Scaling Limits of Prolate Gram Forms, Uniform
Coercivity, and a Trace Formula Toward Weil
Positivity}, preprint, 2026.

\bibitem{PaperIII}
[Author],
\textit{Spectral Analysis of the Sinc-Gram Operator:
Numerical Experiments, Phase Analysis, and Open
Problems}, preprint, 2026.

\bibitem{PaperIV}
[Author],
\textit{Spectral Phase Obstruction for Bandlimited
Operator Families: A No-Go Theorem for
Zeta-Compatible Limits}, preprint, 2026.

\bibitem{PaperVI}
[Author],
\textit{Rigorous Peak-Width Bounds for Prolate
Spheroidal Wave Functions: Uniform Stationary-Phase
Analysis and Multi-Mode Spectral Control},
preprint, 2026.

\end{thebibliography}

\end{document}
